% !TEX root = ../main.tex
\chapter{Fundamentação Teórica}
\label{chap:theoretical_foundation}

Este capítulo apresenta os fundamentos teóricos necessários para a compreensão do trabalho. São abordados os conceitos de armazenamento de dados em DNA, a codificação digital para sequências de nucleotídeos, os tipos de erros e restrições biológicas envolvidos, os códigos de apagamento e \emph{fountain codes}, a formalização dos códigos LT e da distribuição robusta de soliton, o funcionamento do \emph{peeling decoder} e do ripple, a arquitetura do DNA Fountain proposta por Erlich e Zielinski, e as métricas de avaliação utilizadas ao longo dos experimentos.

\section{Armazenamento de Dados em DNA}
\label{sec:dna_data_storage}

A molécula de ácido desoxirribonucleico (DNA) é o biopolímero que atua no armazenamento da informação genética em organismos vivos. Cada fita de DNA é composta por uma sequência linear de nucleotídeos, que são unidades formadas por um grupo fosfato, uma desoxirribose e uma base nitrogenada. As quatro bases nitrogenadas presentes no DNA são adenina (A), citosina (C), guanina (G) e timina (T). Do ponto de vista computacional, a sequência linear de bases pode carregar informação digital arbitrária~\cite{compeau2018}.

Do ponto de vista da teoria da informação, cada nucleotídeo pode representar dois bits de dados digitais, já que existem quatro bases possíveis ($\log_2 4 = 2$). Essa característica confere ao DNA uma densidade informacional teórica de até aproximadamente 455 exabytes por grama~\cite{church2012}, valor expressivamente superior ao de mídias magnéticas ou ópticas.

Além da densidade potencial, o DNA exibe elevada estabilidade físico-química: quando encapsulado sob condições adequadas de baixa temperatura e desidratação (ou protegido em matrizes sintéticas de sílica), o material genético pode preservar cadeias moleculares por milênios com recuperação viável. Diferentemente das mídias magnéticas tradicionais, cuja vida útil típica requer ciclos periódicos de migração a cada 5 a 10 anos, o arquivamento em DNA pode manter os dados estáveis de forma passiva, sem demanda energética contínua para sustentação de campos ou circuitos~\cite{ceze2019}.

Contudo, o armazenamento de dados em DNA também enfrenta desafios significativos. O processo de síntese química de oligonucleotídeos, que consiste na fabricação artificial de sequências curtas de DNA (tipicamente entre 100 e 300 nucleotídeos), apresenta custos elevados e é sujeito a erros. Da mesma forma, o processo de sequenciamento, que realiza a leitura das bases a partir de um pool de moléculas, introduz ruídos e imperfeições na recuperação dos dados. Tecnologias modernas de sequenciamento de nova geração (\emph{Next-Generation Sequencing} --- NGS) e sequenciamento por nanoporos (\emph{Nanopore}) têm reduzido progressivamente esses custos. No entanto, é fundamental distinguir a prova de conceito técnica — já amplamente demonstrada na literatura acadêmica — da viabilidade econômica e operacional em larga escala, que ainda depende da superação de barreiras de custo por gigabyte e da elevada latência de leitura e escrita~\cite{ceze2019}.

O ciclo de vida completo do armazenamento de dados em DNA pode ser descrito em cinco etapas: (1)~codificação digital, na qual o arquivo binário é convertido em sequências de nucleotídeos; (2)~síntese química, na qual as sequências são fabricadas em laboratório; (3)~armazenamento físico do pool de oligonucleotídeos em condições controladas; (4)~sequenciamento, no qual as moléculas são lidas por um dispositivo de sequenciamento; e (5)~decodificação computacional, na qual as leituras obtidas são processadas para reconstruir o arquivo original.

\section{Codificação Digital para DNA}
\label{sec:digital_to_dna_encoding}

A primeira etapa do armazenamento de dados em DNA consiste em converter um arquivo binário em uma representação compatível com as quatro bases nitrogenadas. Para isso, o arquivo de entrada é inicialmente tratado como uma sequência contínua de bytes, que é então dividida em blocos de tamanho fixo $L$ (em bytes). Caso o último bloco não possua exatamente $L$ bytes, é realizado um preenchimento determinístico com bytes nulos (\texttt{0x00}), técnica conhecida como \emph{zero-padding}.

Com os blocos definidos, cada byte é convertido em uma sequência de quatro nucleotídeos por meio de um mapeamento direto de pares de bits para bases. O esquema de mapeamento utilizado neste trabalho é apresentado na Tabela~\ref{tab:mapeamento_bits_dna}.

\begin{table}[htbp]
    \centering
    \caption{Mapeamento de pares de bits para bases nitrogenadas.}
    \label{tab:mapeamento_bits_dna}
    \begin{tabular}{cc}
        \hline
        \textbf{Par de bits} & \textbf{Base} \\
        \hline
        00 & A \\
        01 & C \\
        10 & G \\
        11 & T \\
        \hline
    \end{tabular}
\end{table}

Dessa forma, cada byte (8 bits) é representado por exatamente quatro nucleotídeos, conferindo uma taxa bruta de 2 bits por base para um alfabeto irrestrito de quatro símbolos ($\log_2 4 = 2{,}0$ bits/nt)~\cite{shannon1948}. Na literatura de armazenamento em DNA, distinguem-se com rigor três patamares de capacidade teórica e prática~\cite{erlich2017}:
\begin{enumerate}
    \item \textbf{Capacidade não restrita} ($2{,}0$ bits/nt): limite informacional de Shannon para um canal ideal de quatro símbolos equiprováveis sem qualquer restrição de sequência;
    \item \textbf{Capacidade máxima sob restrições físico-químicas} ($1{,}98$ bits/nt): capacidade combinatória teórica de sequências que respeitam estritamente a janela de conteúdo GC (45\% a 55\%) e o limite de homopolímeros ($\le 3$ bases idênticas consecutivas);
    \item \textbf{Capacidade estimada no modelo de armazenamento de Erlich e Zielinski} ($\approx 1{,}83$ bits/nt): limite operacional de Shannon após descontar a sobrecarga necessária de endereçamento/indexação de sequências desordenadas no pool molecular e a tolerância ao dropout molecular sob as hipóteses do modelo analisado pelos autores.
\end{enumerate}
Nesse contexto, a densidade de informação líquida pioneira de $1{,}57$ bits/nt alcançada experimentalmente por Erlich e Zielinski~\cite{erlich2017} corresponde a $78{,}5\%$ do limite irrestrito de 2 bits/nt, $79{,}3\%$ da capacidade sob restrições físico-químicas ($1{,}98$ bits/nt) e expressivos $85{,}8\%$ (cerca de $86\%$) da capacidade prática estimada no modelo de referência ($\approx 1{,}83$ bits/nt). Este mapeamento de 2 bits por base é biunívoco: a conversão de bytes para DNA e a conversão inversa de DNA para bytes são operações determinísticas e perfeitamente reversíveis na ausência de ruído físico.

A simplicidade deste esquema de mapeamento direto é uma escolha intencional no contexto do DNA Fountain. Como será apresentado nas seções seguintes, a robustez contra erros e perdas não é delegada ao esquema de codificação em si, mas sim à camada de codificação por \emph{fountain codes}, que opera sobre os blocos antes da conversão para DNA.

\section{Erros e Restrições Biológicas}
\label{sec:errors_and_biological_constraints}

Ao longo do ciclo de vida do armazenamento em DNA, duas classes distintas de fenômenos impactam a integridade dos dados: as falhas físicas geradas pelas etapas de síntese e sequenciamento (erros de canal) e as restrições biológicas de aceitabilidade, que definem se uma sequência pode ou não ser sintetizada e amplificada com estabilidade.

\subsection{Erros do Canal}

Os erros do canal ocorrem durante os processos físicos de síntese e sequenciamento, e podem ser classificados em quatro categorias principais:

\begin{itemize}
    \item \textbf{Substituições}: uma base é incorretamente lida ou sintetizada como outra base. Por exemplo, um nucleotídeo G é registrado como A pelo sequenciador.
    
    \item \textbf{Inserções}: uma base adicional é introduzida na sequência durante a síntese ou o sequenciamento, deslocando todas as bases subsequentes.
    
    \item \textbf{Deleções}: uma base é omitida durante a leitura ou a síntese, causando um deslocamento no sentido oposto ao das inserções.
    
    \item \textbf{Apagamentos (\emph{erasures})}: um oligonucleotídeo inteiro é perdido durante o processo, seja por degradação física, cobertura insuficiente no sequenciamento ou falha completa na síntese de uma sequência. Nesse caso, o receptor simplesmente não recebe a sequência, mas sabe que ela está ausente.
\end{itemize}

No contexto deste trabalho, o canal de armazenamento em DNA é modelado de forma simplificada como um \textbf{canal de apagamento} (\emph{erasure channel}), no qual oligonucleotídeos inteiros são perdidos com uma probabilidade $p_{\text{loss}}$, mas as sequências que chegam ao receptor são consideradas íntegras. Essa simplificação é coerente com o escopo delimitado no Capítulo~1 e permite isolar o efeito das perdas sobre a decodificação, sem a complexidade adicional dos erros de substituição, inserção e deleção.

\subsection{Restrições Biológicas de Aceitabilidade}

Independentemente dos erros do canal, nem todas as sequências de DNA podem ser sintetizadas e sequenciadas com a mesma confiabilidade. Certas propriedades da sequência afetam diretamente a eficiência dos processos biológicos, e sequências que violam essas restrições devem ser descartadas antes da síntese. As duas restrições mais relevantes para o DNA Fountain são:

\begin{itemize}
    \item \textbf{Conteúdo GC}: a proporção de bases G e C em relação ao comprimento total da sequência deve estar dentro de uma faixa aceitável, tipicamente entre 40\% e 60\%. Sequências com conteúdo GC muito baixo ou muito alto apresentam temperaturas de desnaturação (\emph{melting temperature}) atípicas, o que dificulta a amplificação por PCR (\emph{Polymerase Chain Reaction}) e introduz viés no sequenciamento~\cite{erlich2017}. Neste trabalho, adotou-se uma janela de GC entre 40\% e 60\%, adotada como configuração do protótipo, mais permissiva que a faixa de 45--55\% utilizada no experimento original.
    
    \item \textbf{Homopolímeros}: sequências que contêm longas repetições consecutivas da mesma base (por exemplo, AAAA ou GGGG) são propensas a erros de síntese e de leitura, especialmente em tecnologias de sequenciamento baseadas em detecção de fluxo. Neste trabalho, o limite máximo adotado é de 3 repetições consecutivas ($\leq 3$), seguindo a configuração proposta por Erlich e Zielinski~\cite{erlich2017}.
\end{itemize}

Essas restrições biológicas \textbf{não constituem erros físicos de canal}, mas sim filtros de aceitabilidade aplicados ainda na etapa de codificação. Quando um droplet gerado viola os limites de GC ou homopolímeros no \emph{screening}, ele é descartado antes do envio e outro droplet é gerado em seu lugar. Essa distinção entre erros estocásticos do canal e critérios de filtragem é essencial para a compreensão do sistema.

\section{Códigos de Apagamento e Fountain Codes}
\label{sec:erasure_codes_and_fountain_codes}

Dado que o canal de DNA é sujeito a perdas de oligonucleotídeos, é necessário utilizar alguma forma de codificação redundante que permita a recuperação dos dados mesmo quando parte das sequências é perdida. Nesse contexto, os \textbf{códigos de apagamento} (\emph{erasure codes}) são a classe de códigos corretores de erros adequada ao problema.

Os códigos de apagamento tradicionais, como os códigos de Reed-Solomon, são códigos de \textbf{taxa fixa} (\emph{fixed-rate codes}). Nesse tipo de código, o codificador gera um número fixo $N$ de pacotes a partir de $K$ blocos originais, com uma taxa de código $R_c = K/N$ congelada em tempo de projeto. Caso a perda real do canal exceda o orçamento estático de paridade ($1 - R_c$), a decodificação torna-se impossível sem uma nova rodada de transmissão ou redimensionamento estrutural do código~\cite{mackay2003}. Em canais de armazenamento molecular em pool, nos quais não há canal de retorno em tempo real (\emph{feedback}) para negociar retransmissões durante a síntese, essa rigidez impõe um dilema entre desperdício excessivo de síntese (se a perda for superestimada) ou perda irreversível de dados (se a perda for subestimada).

Os \textbf{\emph{fountain codes}} (códigos fountain ou códigos de fonte inesgotável), por outro lado, pertencem à classe de códigos de canal para canais de apagamento de \textbf{taxa variável} (\emph{rateless codes}), não devendo ser confundidos com esquemas de codificação de fonte (\emph{source coding}, voltados à compressão). O codificador opera de forma contínua, gerando pacotes codificados (\emph{droplets}) sob demanda até que o receptor colete informação suficiente para reconstruir os $K$ blocos originais. As principais propriedades dessa abordagem são:

\begin{itemize}
    \item O transmissor pode gerar um fluxo potencialmente inesgotável de pacotes codificados a partir dos $K$ blocos originais, flexibilizando a geração de sequências e o dimensionamento da redundância antes da síntese química.
    
    \item O decodificador é capaz de reconstruir a mensagem com alta probabilidade ao coletar uma quantidade moderada de pacotes válidos, tipicamente $N = K(1 + \epsilon)$, com $\epsilon > 0$. Ao contrário de códigos MDS determinísticos (como Reed-Solomon), a recuperação pelo \emph{peeling decoder} é probabilística e depende da conectividade do grafo estocasticamente formado pelos pacotes recebidos.
    
    \item Não há necessidade de canal de retorno (\emph{feedback}) entre receptor e transmissor durante a emissão.
\end{itemize}

Essas propriedades tornam os \emph{fountain codes} vantajosos no planejamento do armazenamento em DNA, no qual o volume exato de oligonucleotídeos que sobreviverão aos processos de síntese, conservação e sequenciamento não pode ser previsto rigidamente. Ressalta-se, contudo, que essa propriedade \emph{rateless} opera na fase de geração digital antes da síntese; uma vez produzida uma biblioteca física fixa de oligonucleotídeos em laboratório, não há criação autônoma de novas moléculas durante a leitura, permanecendo a recuperação condicionada à biblioteca efetivamente armazenada, à profundidade de amostragem no sequenciamento e às perdas moleculares ocorridas.

\section{Códigos LT}
\label{sec:lt_codes}

Os códigos LT (\emph{Luby Transform}), propostos por Michael Luby em 2002~\cite{luby2002}, foram os primeiros \emph{fountain codes} práticos. A ideia central dos códigos LT é gerar cada pacote codificado --- chamado de \textbf{droplet} --- por meio da combinação de um subconjunto aleatório dos blocos originais utilizando a operação de ou-exclusivo (XOR) sobre o corpo finito binário $GF(2)$.

O processo de geração de um droplet segue três passos:

\begin{enumerate}
    \item \textbf{Amostragem do grau}: um valor inteiro $d$ é sorteado de acordo com uma distribuição de probabilidade sobre os graus $\{1, 2, \ldots, K\}$. Esse valor determina quantos blocos fonte serão combinados para formar o droplet.
    
    \item \textbf{Seleção dos blocos}: $d$ blocos distintos são selecionados uniformemente ao acaso dentre os $K$ blocos fonte disponíveis. Os índices dos blocos selecionados são chamados de \textbf{vizinhos} do droplet.
    
    \item \textbf{Combinação XOR}: o payload do droplet é calculado como a soma XOR dos $d$ blocos selecionados:
    
    \begin{equation}
        \label{eq:xor_droplet}
        y = b_{i_1} \oplus b_{i_2} \oplus \cdots \oplus b_{i_d}
    \end{equation}

    onde $b_{i_j}$ denota o $j$-ésimo bloco fonte selecionado e $\oplus$ representa a operação XOR bit a bit.
\end{enumerate}

Esse processo pode ser representado por meio de um \textbf{grafo bipartido}, no qual os $K$ blocos fonte ocupam uma partição (nós de entrada) e os $N$ droplets gerados ocupam a outra partição (nós de saída). Uma aresta conecta um bloco fonte $b_i$ a um droplet $y_j$ se e somente se $b_i$ participou da combinação XOR que gerou $y_j$. O grau de um droplet no grafo corresponde ao número de arestas incidentes, ou seja, ao número de blocos que o compõem.

A escolha da distribuição de probabilidade para o grau $d$ é o fator determinante para o desempenho dos códigos LT. Uma distribuição mal calibrada pode gerar excesso de droplets de alto grau (inúteis nas primeiras iterações da decodificação) ou excesso de droplets de grau 1 (que não contribuem para a diversidade de combinações). A seção seguinte formaliza as distribuições utilizadas.

\section{Distribuição Robusta de Soliton}
\label{sec:robust_soliton_distribution}

A distribuição de probabilidade utilizada para sortear o grau dos droplets é o componente teórico central dos códigos LT. Luby~\cite{luby2002} propôs inicialmente a \textbf{Distribuição Ideal de Soliton} $\rho(d)$, definida como:

\begin{equation}
    \label{eq:ideal_soliton}
    \rho(d) = \begin{cases}
        \dfrac{1}{K}, & \text{se } d = 1 \\[10pt]
        \dfrac{1}{d(d-1)}, & \text{se } d = 2, 3, \ldots, K
    \end{cases}
\end{equation}

A distribuição ideal foi projetada de modo que, em expectativa, exatamente um droplet de grau 1 esteja disponível a cada passo da decodificação. Contudo, na prática, a variância estocástica do processo faz com que o número de droplets de grau 1 disponíveis flutue significativamente, levando com alta probabilidade à interrupção prematura da decodificação.

Para corrigir esse problema, Luby propôs a \textbf{Distribuição Robusta de Soliton} $\mu(d)$, que adiciona à distribuição ideal uma função auxiliar $\tau(d)$ com o objetivo de garantir uma quantidade mínima esperada de droplets de baixo grau. A função $\tau(d)$ é definida como:

\begin{equation}
    \label{eq:tau}
    \tau(d) = \begin{cases}
        \dfrac{R}{d \cdot K}, & \text{se } d = 1, 2, \ldots, \left\lfloor\dfrac{K}{R}\right\rfloor - 1 \\[10pt]
        \dfrac{R \cdot \ln\!\left(\dfrac{R}{\delta}\right)}{K}, & \text{se } d = \left\lfloor\dfrac{K}{R}\right\rfloor \\[10pt]
        0, & \text{se } d > \left\lfloor\dfrac{K}{R}\right\rfloor
    \end{cases}
\end{equation}

onde o parâmetro $R$ e o índice de pico $m$ são definidos por:

\begin{equation}
    \label{eq:R}
    R = c \cdot \ln\!\left(\frac{K}{\delta}\right) \cdot \sqrt{K}, \qquad m = \left\lfloor\frac{K}{R}\right\rfloor
\end{equation}

Na prática computacional, especialmente em regimes com número finito e reduzido de blocos ($K \le 64$), valores extremos de parâmetros podem violar pressupostos analíticos implícitos na formulação assintótica de Luby. Para manter a integridade numérica da amostragem em software, a implementação desenvolvida adota explicitamente duas salvaguardas defensivas de fronteira:
\begin{enumerate}
    \item \textbf{Salvaguarda para $R > K$ ($m \le 0$)}: Quando $c$ é elevado em mensagens curtas, o valor calculado de $R$ pode se aproximar ou superar a própria quantidade de blocos $K$, resultando em $\lfloor K/R \rfloor = 0$. Como os graus dos droplets devem ser inteiros pertencentes a $\{1, \ldots, K\}$, o índice de pico é delimitado computacionalmente por:
    \begin{equation}
        \label{eq:m_adaptado}
        m = \max\left(1, \min\left(K, \left\lfloor\frac{K}{R}\right\rfloor\right)\right)
    \end{equation}
    Na grade avaliada, a correção do índice bruto zero é ativada nas combinações $(K, c, \delta) \in \{(16; 0{,}8; 0{,}10), (16; 0{,}8; 0{,}01), (32; 0{,}8; 0{,}01)\}$. Por exemplo, para $K = 16$, $c = 0{,}8$ e $\delta = 0{,}10$, obtém-se $R \approx 16{,}241$ e $\lfloor K/R \rfloor = 0$, sendo adotado $m = 1$.
    
    \item \textbf{Salvaguarda para $R < \delta$ ($\tau(m) < 0$)}: Caso $R < \delta$, o quociente $R/\delta$ torna-se menor que 1, fazendo com que $\ln(R/\delta) < 0$. Se não tratada, essa eventualidade produziria um termo corretivo negativo em $\tau(m)$, fora das condições usuais da construção, motivo pelo qual a implementação adota uma salvaguarda numérica defensiva:
    \begin{equation}
        \label{eq:tau_adaptado}
        \tau(m) = \max\left(0{,}0, \frac{R}{K} \cdot \ln\!\left(\frac{R}{\delta}\right)\right)
    \end{equation}
    Na grade principal de 216 combinações avaliada neste estudo ($K \ge 16$, $\delta \le 0{,}10$ e $c \ge 0{,}1$), tem-se sempre $R \ge 2{,}03 \gg \delta$, de modo que essa salvaguarda atua como guarda estrutural do software contra configurações hipotéticas degeneradas.
\end{enumerate}

Na implementação computacional, a função corretiva $\tau(d)$ é construída utilizando o índice $m$ após aplicação das salvaguardas das Equações~(\ref{eq:m_adaptado}) e~(\ref{eq:tau_adaptado}), garantindo $1 \le m \le K$ e $\tau(m) \ge 0$ em todos os ramos.

Os parâmetros $c > 0$ e $0 < \delta < 1$ controlam a forma da distribuição robusta:

\begin{itemize}
    \item O parâmetro $c$ escala o valor de $R$. Embora em regimes assintóticos ($K \to \infty$) valores maiores de $c$ aumentem o número esperado de droplets de grau baixo, em regimes com $K$ finito e pequeno o aumento excessivo de $c$ desloca o índice de máximo da função corretiva $\tau(d)$, dado por $m = \lfloor K/R \rfloor$, para graus muito baixos ($m=1$ ou $m=2$). Cabe enfatizar a distinção entre o índice $m$, que maximiza o termo corretivo $\tau(d)$, e a moda global da distribuição $\mu(d)$: para $K = 32$ e $c = 0{,}5$, a translação do pico de $\tau$ para $m=1$ faz com que cerca de 73\% da massa de probabilidade total concentre-se no próprio grau 1, reduzindo expressivamente o grau médio teórico e privando o sistema de conexões intermediárias. Nas configurações avaliadas, o grau médio amostrado diminuiu quando $c$ passou de 0,1 para 0,8, refletindo a redução discreta de $m$; contudo, isso não estabelece monotonicidade estrita em toda a reta real de $c$, haja vista que pequenas oscilações locais podem surgir para variações infinitesimais nos parâmetros da distribuição.
    
    \item O parâmetro $\delta$ orienta o dimensionamento assintótico da cauda da distribuição. De acordo com a formalização original de Luby~\cite{luby2002}, ao receber $K + O(\sqrt{K}\ln^2(K/\delta))$ droplets sob a DRS, o processo de decodificação assegura sucesso na recuperação com probabilidade de pelo menos $1 - \delta$ sob as hipóteses assintóticas do modelo. A garantia teórica assintótica de Luby é apresentada como referência clássica da teoria; ela não se transfere automaticamente às distribuições empíricas de droplets aceitos após screening biológico nem às configurações em blocos finitos ajustadas pelas salvaguardas.
\end{itemize}

A distribuição robusta normalizada $\mu(d)$ é então obtida pela soma de $\rho(d)$ e $\tau(d)$, dividida pela constante de normalização $\beta$:

\begin{equation}
    \label{eq:mu}
    \mu(d) = \frac{\rho(d) + \tau(d)}{\beta}, \quad \text{onde} \quad \beta = \sum_{d=1}^{K} \left[\rho(d) + \tau(d)\right]
\end{equation}

A distribuição resultante $\mu(d)$ satisfaz $\sum_{d=1}^{K} \mu(d) = 1$ e é utilizada para sortear o grau de cada droplet durante a codificação. O grau médio esperado teórico de um droplet sob essa distribuição é formalizado como:

\begin{equation}
    \label{eq:grau_medio}
    \mathbb{E}_\mu[D] = \sum_{d=1}^{K} d \cdot \mu(d)
\end{equation}

Na implementação do simulador, a amostragem de um grau é realizada por meio da função de distribuição acumulada (CDF) de $\mu(d)$: um valor uniforme $u \sim \mathcal{U}(0,1)$ é gerado e o grau $d$ é determinado como o menor valor tal que $\text{CDF}(d) \geq u$. Essa técnica garante que a frequência empírica dos graus sorteados convirja para a distribuição $\mu(d)$ à medida que o número de amostras aumenta.

\section{Decodificador Peeling e Dinâmica do Ripple}
\label{sec:peeling_decoder_and_ripple}

O \textbf{\emph{peeling decoder}} é o algoritmo de decodificação utilizado nos códigos LT para reconstruir os blocos fonte originais a partir dos droplets recebidos. Trata-se de um algoritmo iterativo baseado em propagação de crença (\emph{belief propagation}) sobre o grafo bipartido, que opera exclusivamente com operações XOR no corpo finito binário $GF(2)$.

O princípio de funcionamento do \emph{peeling decoder} pode ser descrito intuitivamente da seguinte forma: se um droplet possui grau 1, ou seja, se ele é a combinação XOR de um único bloco fonte, então seu payload é exatamente igual ao valor desse bloco. Uma vez que o valor de um bloco é descoberto, ele pode ser ``removido'' (por XOR) de todos os demais droplets que o continham, reduzindo o grau desses droplets em uma unidade. Se algum desses droplets atinge grau 1 após a remoção, um novo bloco é descoberto, e o processo se repete.

O algoritmo pode ser formalizado nos seguintes passos:

\begin{enumerate}
    \item \textbf{Inicialização}: Construir o grafo bipartido a partir dos droplets recebidos. Identificar todos os droplets com grau $d = 1$ e inseri-los em uma fila de processamento.
    
    \item \textbf{Iteração}: Enquanto a fila não estiver vazia:
    \begin{enumerate}
        \item Retirar um droplet $y$ de grau 1 da fila.
        \item Identificar o único bloco fonte $b_i$ associado a $y$. Se $b_i$ já foi desvendado em passo anterior (droplet redundante), descartar $y$ e prosseguir para a próxima iteração.
        \item Caso $b_i$ seja inédito, registrar $b_i = y$ como bloco descoberto.
        \item Para cada droplet $y'$ que contém $b_i$ em seus vizinhos:
        \begin{enumerate}
            \item Atualizar o payload de $y'$: $y' \leftarrow y' \oplus b_i$.
            \item Remover $b_i$ da lista de vizinhos de $y'$, reduzindo seu grau em 1.
            \item Se o grau residual de $y'$ for reduzido a zero (equação redundante satisfeita), descartar $y'$. Se o grau de $y'$ se tornar igual a 1, inserir $y'$ na fila de processamento.
        \end{enumerate}
    \end{enumerate}
    
    \item \textbf{Término}: O algoritmo termina quando a fila esvazia. Se todos os $K$ blocos foram descobertos, a decodificação foi bem-sucedida. Caso contrário, o processo de peeling é interrompido sem recuperar a totalidade dos blocos por ausência de equações residuais de grau 1 na fila (o que não exclui a possibilidade de resolução por eliminação gaussiana sobre $GF(2)$, embora o algoritmo de peeling linear tenha estagnado).
\end{enumerate}

O conceito de \textbf{ripple} descreve a dinâmica do algoritmo ao longo das iterações. Na definição original de Luby~\cite{luby2002}, o \emph{ripple} é formalizado como o conjunto de símbolos fonte (variáveis de entrada) cobertos por símbolos codificados liberados, mas ainda não processados pelo decodificador. No presente simulador, a métrica instrumentada corresponde ao comprimento da fila de droplets de grau 1 mantida em memória (\texttt{collections.deque}), a qual armazena operacionalmente os pacotes que atingem grau 1 durante o processo, podendo conter droplets redundantes associados a um mesmo símbolo fonte já desvendado. Essa fila operacional é empregada como indicador prático da dinâmica de decodificação, embora sua cardinalidade não coincida necessariamente com a do ripple teórico original de Luby.

A decodificação atinge sucesso quando todos os $K$ blocos são descobertos antes que a fila se esvazie. Se a fila atinge zero antes que todos os blocos tenham sido recuperados, ocorre a interrupção prematura do peeling (\emph{ripple die-out}). O registro do ripple mínimo como $1{,}0$ no histórico de execução decorre do fato de que a amostragem ocorre no início de cada iteração ativa do laço (\texttt{while ripple:}), de modo que o esvaziamento final (tamanho zero) representa a condição de término do algoritmo.

\section{Arquitetura DNA Fountain}
\label{sec:dna_fountain_architecture}

A arquitetura \textbf{DNA Fountain}, proposta por Erlich e Zielinski em 2017~\cite{erlich2017}, combina os códigos LT com mecanismos específicos para o contexto de armazenamento em DNA, formando um pipeline completo que vai da codificação digital até a reconstrução do arquivo original.

O pipeline do DNA Fountain, conforme modelado neste trabalho, opera da seguinte forma:

\begin{enumerate}
    \item \textbf{Segmentação}: O arquivo binário de entrada é dividido em $K$ blocos de tamanho fixo $L$ bytes. O último bloco recebe \emph{zero-padding} determinístico com bytes nulos caso não possua $L$ bytes completos.
    
    \item \textbf{Geração de droplets}: Para cada droplet, uma \emph{seed} inteira é utilizada como entrada de um PRNG determinístico, que gera o grau $d$ (amostrado de $\mu(d)$) e os $d$ índices dos blocos fonte a serem combinados por XOR.
    
    \item \textbf{Cabeçalho}: Cada droplet recebe um cabeçalho contendo a \emph{seed} codificada. A \emph{seed} atua como o identificador pseudoaleatório compacto que permite regenerar deterministicamente o grau e os índices no decodificador, desde que ambos compartilhem os mesmos parâmetros de protocolo ($K$, algoritmo de PRNG e parametrização da DRS).
    
    \item \textbf{Conversão para DNA}: O conjunto cabeçalho + payload é convertido em uma sequência de nucleotídeos utilizando o mapeamento binário descrito na Seção~\ref{sec:digital_to_dna_encoding}.
    
    \item \textbf{Screening biológico}: A sequência de DNA resultante é avaliada quanto ao conteúdo GC e ao comprimento de homopolímeros. Se a sequência viola alguma das restrições biológicas, o droplet é descartado e um novo droplet é gerado com a próxima \emph{seed} disponível.
    
    \item \textbf{Transmissão pelo canal}: Os oligonucleotídeos que passaram pelo screening são submetidos ao canal de apagamento estocástico com probabilidade de perda $p_{\text{loss}}$.
    
    \item \textbf{Decodificação}: O receptor coleta os oligonucleotídeos recebidos, extrai as \emph{seeds} dos cabeçalhos, regenera os índices dos blocos associados e executa o \emph{peeling decoder} para reconstruir os $K$ blocos fonte originais.
\end{enumerate}

Um aspecto central da arquitetura DNA Fountain é o papel da \emph{seed}. Ao utilizar um PRNG determinístico inicializado com a \emph{seed}, tanto o codificador quanto o decodificador podem calcular independentemente o grau e os índices dos blocos associados a cada droplet. Isso elimina a necessidade de transmitir metadados adicionais sobre a estrutura do grafo bipartido, reduzindo significativamente o overhead de informação lateral.

\section{Métricas de Avaliação}
\label{sec:evaluation_metrics}

Para avaliar quantitativamente o desempenho do simulador sob diferentes configurações, este trabalho utiliza um conjunto de métricas formais que capturam aspectos distintos do processo de codificação e decodificação. As métricas adotadas são:

\begin{itemize}
    \item \textbf{Taxa de sucesso na decodificação} ($P_{\text{succ}}$): a proporção de execuções em que o \emph{peeling decoder} conseguiu recuperar todos os $K$ blocos fonte. Calculada como:
    
    \begin{equation}
        \label{eq:taxa_sucesso}
        P_{\text{succ}} = \frac{\text{número de decodificações completas}}{\text{número total de execuções}}
    \end{equation}
    
    \item \textbf{Overhead relativo nominal e efetivo} ($\epsilon_{\text{tx}}, \epsilon_{\text{rx}}$): a sobrecarga de droplets em relação ao número mínimo teórico $K$. Distinguem-se:
    \begin{enumerate}
        \item \textbf{Overhead nominal de transmissão}: $\epsilon_{\text{tx}} = (N_{\text{tx}} - K) / K$, onde $N_{\text{tx}}$ representa o volume de oligonucleotídeos codificados transmitidos (aprovados no screening biológico);
        \item \textbf{Overhead efetivo recebido}: $\epsilon_{\text{rx}} = (N_{\text{rx}} - K) / K$, onde $N_{\text{rx}}$ representa o total de oligonucleotídeos que chegam ao decodificador após as perdas estocásticas do canal. Em canal ideal ($p_{\text{loss}} = 0{,}0$), tem-se $N_{\text{rx}} = N_{\text{tx}}$ e $\epsilon_{\text{rx}} = \epsilon_{\text{tx}}$.
    \end{enumerate}
    
    \item \textbf{Taxa de rejeição no screening} ($R_{\text{screen}}$): a fração de oligonucleotídeos candidatos gerados que violaram as restrições biológicas físico-químicas (conteúdo GC ou homopolímeros) e foram descartados antes da transmissão:
    
    \begin{equation}
        \label{eq:rejeicao}
        R_{\text{screen}} = \frac{N_{\text{cand}} - N_{\text{tx}}}{N_{\text{cand}}}
    \end{equation}
    
    onde $N_{\text{cand}}$ é o número total de candidatos gerados pelo codificador até atingir a cota de $N_{\text{tx}}$ oligos válidos transmitidos.
    
    \item \textbf{Número de operações XOR}: a contagem acumulada de operações XOR executadas pelo decodificador durante o processo de \emph{peeling}. Cada operação contabilizada corresponde à aplicação da operação XOR bit a bit entre dois blocos completos de $L$ bytes ($\texttt{xor\_blocos}$), executada exatamente uma vez por aresta removida no grafo bipartido durante a propagação de vizinhanças. Essa métrica fornece um indicador direto e agnóstico de plataforma para o esforço algorítmico da decodificação.
    
    \item \textbf{Tamanho do ripple}: registrado ao longo dos passos de decodificação, tanto o valor máximo quanto o valor médio. A trajetória do ripple permite analisar a estabilidade da decodificação e identificar pontos de estagnação.
    
    \item \textbf{Blocos recuperados}: o número de blocos fonte efetivamente reconstruídos ao final da decodificação. Em casos de falha, essa métrica indica o quanto o sistema se aproximou da decodificação completa.
    
    \item \textbf{Tempo de processamento}: medido separadamente para as etapas de codificação e decodificação, em milissegundos. Essa métrica permite avaliar a escalabilidade do simulador em função do número de blocos $K$.
\end{itemize}

Essas métricas são coletadas automaticamente pelo simulador a cada execução e consolidadas em arquivos CSV para análise estatística posterior, conforme detalhado no Capítulo~\ref{chap:experimental_methodology}.
